\documentclass[11pt,a4paper]{article}
\usepackage[margin=1in]{geometry}
\usepackage{amsmath,booktabs,hyperref,float}

\title{Iteration 027: PLS Regression Init + FIR}
\author{Autoresearch Loop}
\date{April 8, 2026}

\begin{document}
\maketitle

\begin{abstract}
PLS regression initialization (covariance-maximizing) followed by FIR fine-tuning
achieves $r=0.377$, comparable to the CF-initialized FIR ($r=0.378$). The
PLS and OLS solutions are nearly identical for this data, confirming that
MSE-optimal and covariance-optimal solutions coincide.
\end{abstract}

\section{Method}
Flatten training windows to $(N \cdot T, C_{in})$ and $(N \cdot T, C_{out})$.
Fit PLS regression with 12 components, extract spatial filter $W$, initialize
FIR center-tap. Fine-tune 150 epochs with combined loss + corr validation.

\section{Results}
\begin{table}[H]
\centering
\begin{tabular}{lrrr}
\toprule
Model & Mean $r$ & Std $r$ & SNR (dB) \\
\midrule
CF + FIR (iter017) & 0.378 & 0.078 & 0.61 \\
PLS + FIR & 0.377 & 0.077 & 0.61 \\
\bottomrule
\end{tabular}
\end{table}

\section{Analysis}
PLS and OLS give essentially the same spatial filter because the cross-covariance
matrix $R_{YX}$ has relatively low rank (12 in-ear channels). Both find the
same dominant directions. The initialization matters less than the fine-tuning.

\end{document}
